\section{Recipe 实践清单：如何把数据变成能力}

Data Recipe 不是口号，而是一套实验流程。一个团队如果说自己有数据优势，至少要能回答：数据从哪来，如何筛选，如何混合，如何进入训练，如何评估，失败如何回流，下一轮如何调整。

\subsection{七步实践清单}

\noindent\begin{longtable}{p{0.16\textwidth}p{0.36\textwidth}p{0.39\textwidth}}
\toprule
步骤 & 操作 & 检查点 \\
\midrule
1. 定义能力 & 明确要提升什么任务能力 & 是感知、规划、动作、恢复，还是泛化？ \\
2. 找数据源 & 真实、仿真、合成、互联网、人类反馈 & 数据是否覆盖目标能力？ \\
3. 清洗过滤 & 去重、去噪、去无关、去污染 & 是否误删长尾和失败案例？ \\
4. 设配比 & 确定不同来源和难度比例 & 是否造成能力偏科？ \\
5. 训练课程 & 决定先学什么后学什么 & 是否符合能力递进？ \\
6. 评估反馈 & 用 benchmark/仿真/真实任务检查 & 指标是否和真实部署相关？ \\
7. 闭环回流 & 把失败、长尾和新任务回到数据池 & 下一轮数据是否更有边际收益？ \\
\bottomrule
\end{longtable}

\begin{knowledgebox}{Recipe 的可复制与不可复制}
流程可以复制，但具体配方难复制。因为配方来自模型、任务、硬件、评测和团队经验的交互；换一个目标能力，配方就可能要重做。
\end{knowledgebox}

\subsection{本章小结}

把数据变成能力，需要从目标能力出发，而不是从手里有什么数据出发。Recipe 是能力导向的数据工程。

